% =============================================
% ChronoVault — Rapport de Projet
% Auteurs : ELHARCHY Mostafa, EL MESKI Karim
% =============================================

\documentclass[12pt,a4paper]{report}

% ── Packages ─────────────────────────────────
\usepackage[utf8]{inputenc}
\usepackage[T1]{fontenc}
\usepackage[french]{babel}
\usepackage{geometry}
\usepackage{graphicx}
\usepackage{hyperref}
\usepackage{xcolor}
\usepackage{listings}
\usepackage{tikz}
\usepackage{fancyhdr}
\usepackage{titlesec}
\usepackage{enumitem}
\usepackage{booktabs}
\usepackage{tabularx}
\usepackage{caption}
\usepackage{float}
\usepackage{tcolorbox}
\usepackage{multicol}

% ── Géométrie ────────────────────────────────
\geometry{margin=2.5cm}

% ── Couleurs ─────────────────────────────────
\definecolor{cvgreen}{HTML}{16A34A}
\definecolor{cvdark}{HTML}{171717}
\definecolor{cvgray}{HTML}{737373}
\definecolor{codebg}{HTML}{F5F5F5}
\definecolor{codeframe}{HTML}{D4D4D4}

% ── Listings (code) ─────────────────────────
\lstset{
    basicstyle=\ttfamily\footnotesize,
    backgroundcolor=\color{codebg},
    frame=single,
    rulecolor=\color{codeframe},
    breaklines=true,
    breakatwhitespace=true,
    tabsize=4,
    showstringspaces=false,
    captionpos=b,
    numbers=left,
    numberstyle=\tiny\color{cvgray},
    keywordstyle=\color{cvgreen}\bfseries,
    commentstyle=\color{cvgray}\itshape,
    stringstyle=\color{cvdark},
    xleftmargin=1em,
    framexleftmargin=1em
}

% ── En-têtes / pieds de page ─────────────────
\pagestyle{fancy}
\fancyhf{}
\fancyhead[L]{\small\textcolor{cvgray}{ChronoVault}}
\fancyhead[R]{\small\textcolor{cvgray}{\leftmark}}
\fancyfoot[C]{\thepage}
\renewcommand{\headrulewidth}{0.4pt}

% ── Chapitres stylisés ──────────────────────
\titleformat{\chapter}[display]
{\normalfont\huge\bfseries\color{cvdark}}
{\textcolor{cvgreen}{\chaptertitlename\ \thechapter}}{20pt}{\Huge}

% ── Liens ────────────────────────────────────
\hypersetup{
    colorlinks=true,
    linkcolor=cvgreen,
    urlcolor=cvgreen,
    citecolor=cvgreen
}

% ══════════════════════════════════════════════
%  PAGE DE GARDE
% ══════════════════════════════════════════════
\begin{document}

\begin{titlepage}
    \centering
    \vspace*{2cm}

    {\Huge\bfseries\textcolor{cvdark}{ChronoVault}\par}
    \vspace{0.5cm}
    {\Large\textcolor{cvgreen}{Coffre-fort de souvenirs collaboratif}\par}
    \vspace{1cm}

    \rule{0.6\textwidth}{1pt}\par
    \vspace{1.5cm}

    {\large\textbf{Rapport de Projet}\par}
    \vspace{2cm}

    {\large
        \textbf{Réalisé par :}\par
        \vspace{0.5cm}
        ELHARCHY Mostafa \par
        EL MESKI Karim \par
    }

    \vspace{2cm}

    {\large
        \textbf{Technologies utilisées :}\par
        \vspace{0.3cm}
        PHP 8 \textbullet{} MySQL/MariaDB \textbullet{} HTML5 \textbullet{} Tailwind CSS \textbullet{} JavaScript\par
    }

    \vfill
    {\large Année universitaire 2025 -- 2026\par}
\end{titlepage}

% ── Table des matières ───────────────────────
\tableofcontents
\newpage

% ══════════════════════════════════════════════
%  CHAPITRE 1 : INTRODUCTION
% ══════════════════════════════════════════════
\chapter{Introduction}

\section{Contexte du projet}

À l'ère du numérique, les souvenirs sont de plus en plus partagés sous forme de photos, de messages et de vidéos. Cependant, la notion de \textit{capsule temporelle} --- un conteneur que l'on remplit de souvenirs et que l'on n'ouvre qu'à une date future --- reste un concept captivant qui n'a pas encore été pleinement exploité dans le monde des applications web.

Le projet \textbf{ChronoVault} s'inscrit dans cette optique : il propose une \textbf{application web collaborative} permettant à plusieurs utilisateurs de créer des capsules temporelles numériques, d'y déposer des souvenirs (textes, images), et de ne les révéler qu'à une date d'ouverture définie à l'avance.

\section{Problématique}

Comment concevoir une application web sécurisée et intuitive permettant à un groupe d'utilisateurs de :
\begin{itemize}[noitemsep]
    \item Créer et gérer des capsules temporelles collaboratives ?
    \item Inviter d'autres utilisateurs à participer à une capsule ?
    \item Ajouter des souvenirs (textes et images) dans une capsule verrouillée ?
    \item Garantir que le contenu ne soit accessible qu'après la date d'ouverture ?
\end{itemize}

\section{Objectifs}

Les objectifs principaux de ce projet sont :
\begin{enumerate}[noitemsep]
    \item \textbf{Authentification sécurisée} : inscription, connexion et gestion de sessions avec hachage bcrypt.
    \item \textbf{Gestion des capsules} : création, consultation, suppression et verrouillage temporel.
    \item \textbf{Collaboration} : système d'invitations permettant d'ajouter des membres à une capsule.
    \item \textbf{Souvenirs multimédia} : ajout de messages texte et d'images (upload sécurisé).
    \item \textbf{Interface moderne} : design responsive avec mode sombre/clair.
    \item \textbf{Architecture SPA} : navigation fluide sans rechargement de page.
\end{enumerate}

\section{Technologies utilisées}

\begin{table}[H]
    \centering
    \begin{tabularx}{\textwidth}{l X}
        \toprule
        \textbf{Couche} & \textbf{Technologies} \\
        \midrule
        Back-end & PHP 8, API REST, sessions PHP, bcrypt \\
        Base de données & MySQL / MariaDB, PDO \\
        Front-end & HTML5, JavaScript (Vanilla), Tailwind CSS (CDN) \\
        Serveur & Apache (XAMPP), .htaccess pour le routage SPA \\
        Outils & Git, VS Code, XAMPP \\
        \bottomrule
    \end{tabularx}
    \caption{Stack technologique du projet}
\end{table}


% ══════════════════════════════════════════════
%  CHAPITRE 2 : CONCEPTION
% ══════════════════════════════════════════════
\chapter{Conception}

Ce chapitre présente la modélisation UML du projet ChronoVault à travers trois types de diagrammes : cas d'utilisation, classes et séquence.

% ─────────────────────────────────────────────
\section{Diagramme de cas d'utilisation}

Le diagramme de cas d'utilisation identifie les acteurs du système et les fonctionnalités auxquelles ils ont accès.

\subsection{Acteurs}

\begin{itemize}[noitemsep]
    \item \textbf{Visiteur} : utilisateur non authentifié, il peut uniquement s'inscrire ou se connecter.
    \item \textbf{Utilisateur authentifié} : peut créer des capsules, inviter des membres, ajouter des souvenirs et gérer ses invitations.
\end{itemize}

\subsection{Cas d'utilisation}

\begin{figure}[H]
    \centering
    \begin{tcolorbox}[colback=codebg, colframe=codeframe, width=\textwidth, title={\textbf{Diagramme de cas d'utilisation — ChronoVault}}]
        \small
        \textbf{Acteur : Visiteur}
        \begin{itemize}[noitemsep, leftmargin=1.5em]
            \item S'inscrire
            \item Se connecter
        \end{itemize}

        \vspace{0.3cm}
        \textbf{Acteur : Utilisateur authentifié}
        \begin{itemize}[noitemsep, leftmargin=1.5em]
            \item Se déconnecter
            \item Créer une capsule temporelle
            \item Consulter ses capsules
            \item Supprimer une capsule (créateur uniquement)
            \item Inviter un utilisateur à une capsule
            \item Accepter / Refuser une invitation
            \item Ajouter un souvenir (texte ou image)
            \item Supprimer un souvenir
            \item Consulter les souvenirs (si capsule ouverte)
            \item Basculer mode sombre / clair
        \end{itemize}
    \end{tcolorbox}
    \caption{Cas d'utilisation du système ChronoVault}
\end{figure}

\noindent
\textbf{Description textuelle UML :}

\begin{verbatim}
@startuml
left to right direction

actor "Visiteur" as V
actor "Utilisateur" as U

rectangle ChronoVault {
    V --> (S'inscrire)
    V --> (Se connecter)

    U --> (Se déconnecter)
    U --> (Créer une capsule)
    U --> (Consulter ses capsules)
    U --> (Supprimer une capsule)
    U --> (Inviter un utilisateur)
    U --> (Accepter / Refuser invitation)
    U --> (Ajouter un souvenir)
    U --> (Supprimer un souvenir)
    U --> (Consulter les souvenirs)
    U --> (Basculer thème)

    (S'inscrire) .> (Se connecter) : <<include>>
}
@enduml
\end{verbatim}


% ─────────────────────────────────────────────
\section{Diagramme de classes}

Le diagramme de classes modélise la structure des données et les relations entre les entités du système.

\begin{figure}[H]
    \centering
    \begin{tcolorbox}[colback=codebg, colframe=codeframe, width=\textwidth, title={\textbf{Diagramme de classes — ChronoVault}}]
        \small
        \textbf{User}
        \begin{itemize}[noitemsep, leftmargin=1.5em]
            \item \texttt{- id : INT (PK, AUTO\_INCREMENT)}
            \item \texttt{- username : VARCHAR(50) UNIQUE}
            \item \texttt{- email : VARCHAR(100) UNIQUE}
            \item \texttt{- password\_hash : VARCHAR(255)}
            \item \texttt{- avatar : VARCHAR(255)}
            \item \texttt{- created\_at : DATETIME}
            \item \texttt{- updated\_at : DATETIME}
        \end{itemize}

        \vspace{0.4cm}
        \textbf{Capsule}
        \begin{itemize}[noitemsep, leftmargin=1.5em]
            \item \texttt{- id : INT (PK, AUTO\_INCREMENT)}
            \item \texttt{- title : VARCHAR(150)}
            \item \texttt{- description : TEXT}
            \item \texttt{- opening\_date : DATE}
            \item \texttt{- creator\_id : INT (FK → User.id)}
            \item \texttt{- is\_opened : TINYINT(1)}
            \item \texttt{- cover\_image : VARCHAR(255)}
            \item \texttt{- created\_at : DATETIME}
            \item \texttt{- updated\_at : DATETIME}
        \end{itemize}

        \vspace{0.4cm}
        \textbf{CapsuleMember}
        \begin{itemize}[noitemsep, leftmargin=1.5em]
            \item \texttt{- id : INT (PK, AUTO\_INCREMENT)}
            \item \texttt{- capsule\_id : INT (FK → Capsule.id)}
            \item \texttt{- user\_id : INT (FK → User.id)}
            \item \texttt{- role : ENUM('creator', 'member')}
            \item \texttt{- joined\_at : DATETIME}
        \end{itemize}

        \vspace{0.4cm}
        \textbf{Invitation}
        \begin{itemize}[noitemsep, leftmargin=1.5em]
            \item \texttt{- id : INT (PK, AUTO\_INCREMENT)}
            \item \texttt{- capsule\_id : INT (FK → Capsule.id)}
            \item \texttt{- invited\_by : INT (FK → User.id)}
            \item \texttt{- invited\_user\_id : INT (FK → User.id)}
            \item \texttt{- status : ENUM('pending', 'accepted', 'declined')}
            \item \texttt{- created\_at : DATETIME}
        \end{itemize}

        \vspace{0.4cm}
        \textbf{Memory}
        \begin{itemize}[noitemsep, leftmargin=1.5em]
            \item \texttt{- id : INT (PK, AUTO\_INCREMENT)}
            \item \texttt{- capsule\_id : INT (FK → Capsule.id)}
            \item \texttt{- user\_id : INT (FK → User.id)}
            \item \texttt{- text\_content : TEXT}
            \item \texttt{- image\_path : VARCHAR(255)}
            \item \texttt{- created\_at : DATETIME}
            \item \texttt{- updated\_at : DATETIME}
        \end{itemize}

        \vspace{0.5cm}
        \textbf{Relations :}
        \begin{itemize}[noitemsep, leftmargin=1.5em]
            \item User \texttt{1} --- \texttt{0..*} Capsule \hfill (un utilisateur crée 0 à N capsules)
            \item User \texttt{1} --- \texttt{0..*} CapsuleMember \hfill (un utilisateur est membre de 0 à N capsules)
            \item Capsule \texttt{1} --- \texttt{1..*} CapsuleMember \hfill (une capsule a au moins 1 membre)
            \item User \texttt{1} --- \texttt{0..*} Invitation \hfill (un utilisateur envoie/reçoit des invitations)
            \item Capsule \texttt{1} --- \texttt{0..*} Invitation \hfill (une capsule peut avoir des invitations)
            \item User \texttt{1} --- \texttt{0..*} Memory \hfill (un utilisateur crée des souvenirs)
            \item Capsule \texttt{1} --- \texttt{0..*} Memory \hfill (une capsule contient des souvenirs)
        \end{itemize}
    \end{tcolorbox}
    \caption{Diagramme de classes du système ChronoVault}
\end{figure}

\noindent
\textbf{Description textuelle UML (PlantUML) :}

\begin{verbatim}
@startuml
class User {
    - id : int
    - username : string
    - email : string
    - password_hash : string
    - avatar : string
    - created_at : datetime
    - updated_at : datetime
}

class Capsule {
    - id : int
    - title : string
    - description : string
    - opening_date : date
    - creator_id : int
    - is_opened : boolean
    - cover_image : string
    - created_at : datetime
}

class CapsuleMember {
    - id : int
    - capsule_id : int
    - user_id : int
    - role : enum
    - joined_at : datetime
}

class Invitation {
    - id : int
    - capsule_id : int
    - invited_by : int
    - invited_user_id : int
    - status : enum
    - created_at : datetime
}

class Memory {
    - id : int
    - capsule_id : int
    - user_id : int
    - text_content : string
    - image_path : string
    - created_at : datetime
}

User "1" -- "0..*" Capsule : crée
User "1" -- "0..*" CapsuleMember : participe
Capsule "1" -- "1..*" CapsuleMember : contient
User "1" -- "0..*" Invitation : envoie/reçoit
Capsule "1" -- "0..*" Invitation : a
User "1" -- "0..*" Memory : rédige
Capsule "1" -- "0..*" Memory : contient
@enduml
\end{verbatim}


% ─────────────────────────────────────────────
\section{Diagrammes de séquence}

Les diagrammes de séquence illustrent les interactions entre l'utilisateur, le front-end (SPA) et le back-end (API PHP) pour les scénarios clés.

\subsection{Séquence : Inscription}

\begin{figure}[H]
    \centering
    \begin{tcolorbox}[colback=codebg, colframe=codeframe, width=\textwidth, title={\textbf{Diagramme de séquence — Inscription}}]
        \small
        \begin{enumerate}[noitemsep]
            \item \textbf{Utilisateur} $\rightarrow$ \textbf{SPA} : Remplit le formulaire (pseudo, email, mot de passe)
            \item \textbf{SPA} $\rightarrow$ \textbf{API} : \texttt{POST /api/auth.php?action=register} \{username, email, password\}
            \item \textbf{API} $\rightarrow$ \textbf{BDD} : Vérifier unicité (username, email)
            \item \textbf{BDD} $\rightarrow$ \textbf{API} : Résultat de la vérification
            \item \textbf{API} $\rightarrow$ \textbf{BDD} : \texttt{INSERT INTO users} (password haché bcrypt)
            \item \textbf{API} $\rightarrow$ \textbf{API} : Créer session PHP (\texttt{\$\_SESSION['user\_id']})
            \item \textbf{API} $\rightarrow$ \textbf{SPA} : \texttt{\{success: true, user: \{...\}\}}
            \item \textbf{SPA} $\rightarrow$ \textbf{Utilisateur} : Toast « Compte créé ! » + redirection Dashboard
        \end{enumerate}
    \end{tcolorbox}
    \caption{Séquence d'inscription d'un utilisateur}
\end{figure}

\noindent
\textbf{PlantUML :}

\begin{verbatim}
@startuml
actor Utilisateur
participant "SPA (JavaScript)" as SPA
participant "API PHP" as API
database "MySQL" as DB

Utilisateur -> SPA : Remplit formulaire inscription
SPA -> API : POST /api/auth.php?action=register
API -> DB : SELECT (vérifier unicité)
DB --> API : OK / Doublon
alt Utilisateur unique
    API -> DB : INSERT INTO users
    DB --> API : userId
    API -> API : $_SESSION['user_id'] = userId
    API --> SPA : {success: true, user: {...}}
    SPA --> Utilisateur : Toast + Dashboard
else Doublon
    API --> SPA : {success: false, error: "..."}
    SPA --> Utilisateur : Toast erreur
end
@enduml
\end{verbatim}


\subsection{Séquence : Création d'une capsule}

\begin{figure}[H]
    \centering
    \begin{tcolorbox}[colback=codebg, colframe=codeframe, width=\textwidth, title={\textbf{Diagramme de séquence — Création d'une capsule}}]
        \small
        \begin{enumerate}[noitemsep]
            \item \textbf{Utilisateur} $\rightarrow$ \textbf{SPA} : Remplit titre, description, date d'ouverture
            \item \textbf{SPA} $\rightarrow$ \textbf{API} : \texttt{POST /api/capsules/?action=create} \{title, description, opening\_date\}
            \item \textbf{API} $\rightarrow$ \textbf{API} : Vérifier authentification (\texttt{requireAuth()})
            \item \textbf{API} $\rightarrow$ \textbf{BDD} : \texttt{INSERT INTO capsules}
            \item \textbf{BDD} $\rightarrow$ \textbf{API} : capsuleId
            \item \textbf{API} $\rightarrow$ \textbf{BDD} : \texttt{INSERT INTO capsule\_members} (rôle = 'creator')
            \item \textbf{API} $\rightarrow$ \textbf{SPA} : \texttt{\{success: true, capsule\_id: ...\}}
            \item \textbf{SPA} $\rightarrow$ \textbf{Utilisateur} : Toast « Capsule créée ! » + redirection Dashboard
        \end{enumerate}
    \end{tcolorbox}
    \caption{Séquence de création d'une capsule temporelle}
\end{figure}


\subsection{Séquence : Ajout d'un souvenir}

\begin{figure}[H]
    \centering
    \begin{tcolorbox}[colback=codebg, colframe=codeframe, width=\textwidth, title={\textbf{Diagramme de séquence — Ajout d'un souvenir}}]
        \small
        \begin{enumerate}[noitemsep]
            \item \textbf{Utilisateur} $\rightarrow$ \textbf{SPA} : Rédige un texte et/ou sélectionne une image
            \item \textbf{SPA} $\rightarrow$ \textbf{API} : \texttt{POST /api/memories/?action=add} (FormData : capsule\_id, text\_content, image)
            \item \textbf{API} $\rightarrow$ \textbf{API} : Vérifier authentification + appartenance à la capsule
            \item \textbf{API} $\rightarrow$ \textbf{API} : Valider l'image (type, taille $\leq$ 5 Mo)
            \item \textbf{API} $\rightarrow$ \textbf{Serveur} : Sauvegarder l'image dans \texttt{/uploads/}
            \item \textbf{API} $\rightarrow$ \textbf{BDD} : \texttt{INSERT INTO memories}
            \item \textbf{API} $\rightarrow$ \textbf{SPA} : \texttt{\{success: true, memory: \{...\}\}}
            \item \textbf{SPA} $\rightarrow$ \textbf{Utilisateur} : Toast « Souvenir ajouté ! » + rafraîchissement
        \end{enumerate}
    \end{tcolorbox}
    \caption{Séquence d'ajout d'un souvenir dans une capsule}
\end{figure}


\subsection{Séquence : Invitation d'un utilisateur}

\begin{figure}[H]
    \centering
    \begin{tcolorbox}[colback=codebg, colframe=codeframe, width=\textwidth, title={\textbf{Diagramme de séquence — Invitation}}]
        \small
        \begin{enumerate}[noitemsep]
            \item \textbf{Créateur} $\rightarrow$ \textbf{SPA} : Saisit le pseudo de l'utilisateur à inviter
            \item \textbf{SPA} $\rightarrow$ \textbf{API} : \texttt{POST /api/invitations/?action=send} \{capsule\_id, username\}
            \item \textbf{API} $\rightarrow$ \textbf{BDD} : Vérifier que l'utilisateur existe
            \item \textbf{API} $\rightarrow$ \textbf{BDD} : Vérifier qu'il n'est pas déjà membre / invité
            \item \textbf{API} $\rightarrow$ \textbf{BDD} : \texttt{INSERT INTO invitations} (status = 'pending')
            \item \textbf{API} $\rightarrow$ \textbf{SPA} : \texttt{\{success: true\}}
            \item \textbf{SPA} $\rightarrow$ \textbf{Créateur} : Toast « Invitation envoyée ! »
        \end{enumerate}

        \vspace{0.3cm}
        \textbf{Côté invité :}
        \begin{enumerate}[noitemsep]
            \item \textbf{Invité} $\rightarrow$ \textbf{SPA} : Consulte ses invitations
            \item \textbf{SPA} $\rightarrow$ \textbf{API} : \texttt{GET /api/invitations/?action=list}
            \item \textbf{Invité} $\rightarrow$ \textbf{SPA} : Clique « Accepter »
            \item \textbf{SPA} $\rightarrow$ \textbf{API} : \texttt{POST /api/invitations/?action=accept}
            \item \textbf{API} $\rightarrow$ \textbf{BDD} : \texttt{UPDATE invitations SET status='accepted'}
            \item \textbf{API} $\rightarrow$ \textbf{BDD} : \texttt{INSERT INTO capsule\_members}
        \end{enumerate}
    \end{tcolorbox}
    \caption{Séquence d'invitation et d'acceptation}
\end{figure}


% ══════════════════════════════════════════════
%  CHAPITRE 3 : RÉALISATION
% ══════════════════════════════════════════════
\chapter{Réalisation}

Ce chapitre décrit l'implémentation technique du projet ChronoVault, l'architecture des fichiers, les composants principaux et les captures d'écran de l'application.

\section{Architecture du projet}

Le projet suit une architecture \textbf{SPA (Single Page Application)} avec une API REST côté serveur.

\begin{figure}[H]
    \centering
    \begin{tcolorbox}[colback=codebg, colframe=codeframe, width=0.75\textwidth, title={\textbf{Arborescence du projet}}]
        \small
\begin{verbatim}
chronovault/
├── index.html              (Point d'entrée SPA)
├── logo.png                (Logo de l'application)
├── .htaccess               (Routage Apache SPA)
├── js/
│   ├── app.js              (Logique SPA principale)
│   └── api.js              (Client HTTP pour l'API)
├── api/
│   ├── auth.php            (Authentification)
│   ├── config/
│   │   └── database.php    (Connexion PDO)
│   ├── includes/
│   │   └── helpers.php     (Fonctions utilitaires)
│   ├── capsules/
│   │   └── index.php       (CRUD capsules)
│   ├── invitations/
│   │   └── index.php       (Gestion invitations)
│   └── memories/
│       └── index.php       (Gestion souvenirs)
├── database/
│   └── chronovault.sql     (Schéma SQL)
└── uploads/                (Images uploadées)
\end{verbatim}
    \end{tcolorbox}
    \caption{Structure des fichiers du projet}
\end{figure}


\section{Base de données}

La base de données \texttt{chronovault} utilise \textbf{MariaDB} avec le moteur \textbf{InnoDB} pour garantir l'intégrité référentielle via les clés étrangères.

\subsection{Schéma relationnel}

\begin{table}[H]
    \centering
    \begin{tabularx}{\textwidth}{l l X}
        \toprule
        \textbf{Table} & \textbf{Clé primaire} & \textbf{Description} \\
        \midrule
        \texttt{users} & \texttt{id} & Utilisateurs (pseudo unique, email unique, mot de passe haché) \\
        \texttt{capsules} & \texttt{id} & Capsules temporelles (titre, description, date d'ouverture) \\
        \texttt{capsule\_members} & \texttt{id} & Relation N:N entre utilisateurs et capsules (rôle) \\
        \texttt{invitations} & \texttt{id} & Invitations en attente, acceptées ou refusées \\
        \texttt{memories} & \texttt{id} & Souvenirs : texte et/ou chemin d'image \\
        \bottomrule
    \end{tabularx}
    \caption{Tables de la base de données}
\end{table}

\subsection{Extrait du schéma SQL}

\begin{lstlisting}[language=SQL, caption={Création de la table capsules}]
CREATE TABLE capsules (
    id INT AUTO_INCREMENT PRIMARY KEY,
    title VARCHAR(150) NOT NULL,
    description TEXT,
    opening_date DATE NOT NULL,
    creator_id INT NOT NULL,
    is_opened TINYINT(1) DEFAULT 0,
    cover_image VARCHAR(255) DEFAULT NULL,
    created_at DATETIME DEFAULT CURRENT_TIMESTAMP,
    updated_at DATETIME DEFAULT CURRENT_TIMESTAMP
        ON UPDATE CURRENT_TIMESTAMP,
    FOREIGN KEY (creator_id)
        REFERENCES users(id) ON DELETE CASCADE
) ENGINE=InnoDB;
\end{lstlisting}

\subsection{Index pour l'optimisation}

\begin{lstlisting}[language=SQL, caption={Index créés pour accélérer les requêtes}]
CREATE INDEX idx_capsules_opening_date ON capsules(opening_date);
CREATE INDEX idx_capsules_creator ON capsules(creator_id);
CREATE INDEX idx_memories_capsule ON memories(capsule_id);
CREATE INDEX idx_memories_user ON memories(user_id);
CREATE INDEX idx_invitations_user ON invitations(invited_user_id);
CREATE INDEX idx_invitations_status ON invitations(status);
\end{lstlisting}


\section{Back-end : API REST en PHP}

L'API REST est construite en \textbf{PHP 8} natif, sans framework, avec \textbf{PDO} pour l'accès à la base de données.

\subsection{Points d'entrée de l'API}

\begin{table}[H]
    \centering
    \begin{tabularx}{\textwidth}{l l l X}
        \toprule
        \textbf{Méthode} & \textbf{Endpoint} & \textbf{Action} & \textbf{Description} \\
        \midrule
        POST & \texttt{/api/auth.php} & register & Inscription \\
        POST & \texttt{/api/auth.php} & login & Connexion \\
        GET  & \texttt{/api/auth.php} & logout & Déconnexion \\
        GET  & \texttt{/api/auth.php} & me & Utilisateur courant \\
        \midrule
        GET  & \texttt{/api/capsules/} & list & Lister mes capsules \\
        GET  & \texttt{/api/capsules/} & detail & Détail d'une capsule \\
        POST & \texttt{/api/capsules/} & create & Créer une capsule \\
        POST & \texttt{/api/capsules/} & delete & Supprimer une capsule \\
        \midrule
        GET  & \texttt{/api/invitations/} & list & Mes invitations \\
        POST & \texttt{/api/invitations/} & send & Envoyer une invitation \\
        POST & \texttt{/api/invitations/} & accept & Accepter \\
        POST & \texttt{/api/invitations/} & decline & Refuser \\
        \midrule
        POST & \texttt{/api/memories/} & add & Ajouter un souvenir \\
        POST & \texttt{/api/memories/} & delete & Supprimer un souvenir \\
        \bottomrule
    \end{tabularx}
    \caption{Endpoints de l'API REST}
\end{table}

\subsection{Sécurité}

\begin{itemize}[noitemsep]
    \item \textbf{Hachage des mots de passe} : utilisation de \texttt{password\_hash()} avec l'algorithme \texttt{PASSWORD\_BCRYPT}.
    \item \textbf{Sessions PHP} : l'identifiant utilisateur est stocké dans \texttt{\$\_SESSION['user\_id']} côté serveur.
    \item \textbf{Validation des entrées} : vérification de la longueur, du format email, de l'unicité en base.
    \item \textbf{Upload sécurisé} : vérification du type MIME (JPEG, PNG, GIF, WEBP), taille max 5 Mo, nom de fichier unique généré par \texttt{uniqid()}.
    \item \textbf{Requêtes préparées} : toutes les requêtes SQL utilisent des \textbf{prepared statements} PDO pour prévenir les injections SQL.
    \item \textbf{Gestion d'erreurs} : gestionnaire global d'exceptions renvoyant du JSON structuré.
\end{itemize}

\subsection{Exemple de code : authentification}

\begin{lstlisting}[language=PHP, caption={Extrait de auth.php — Inscription}]
case 'register':
    $data = getJsonBody();
    $username = trim($data['username'] ?? '');
    $email    = trim($data['email'] ?? '');
    $password = $data['password'] ?? '';

    // Validations
    if (!$username || !$email || !$password)
        jsonError('Tous les champs sont obligatoires.');

    // Verifier unicite
    $stmt = $db->prepare(
        'SELECT id FROM users WHERE username = ? OR email = ?'
    );
    $stmt->execute([$username, $email]);
    if ($stmt->fetch())
        jsonError('Ce pseudo ou email est deja utilise.');

    // Hachage + insertion
    $hash = password_hash($password, PASSWORD_BCRYPT);
    $stmt = $db->prepare(
        'INSERT INTO users (username, email, password_hash)
         VALUES (?, ?, ?)'
    );
    $stmt->execute([$username, $email, $hash]);
\end{lstlisting}


\section{Front-end : SPA en JavaScript}

Le front-end est une \textbf{Single Page Application} (SPA) écrite en JavaScript vanilla, sans framework. La mise en forme utilise \textbf{Tailwind CSS} via CDN.

\subsection{Principe de fonctionnement}

L'objet \texttt{app} centralise toute la logique de l'application :

\begin{itemize}[noitemsep]
    \item \textbf{Navigation} : la méthode \texttt{navigate(page, params)} vide le conteneur \texttt{\#app} et appelle la méthode de rendu correspondante (\texttt{renderLogin}, \texttt{renderDashboard}, etc.).
    \item \textbf{Rendu dynamique} : chaque page est générée par du HTML injecté via \texttt{innerHTML}, avec des \texttt{addEventListener} pour gérer les interactions.
    \item \textbf{Client API} : l'objet \texttt{api} (dans \texttt{api.js}) encapsule tous les appels \texttt{fetch()} vers l'API REST.
    \item \textbf{Notifications} : un système de \textit{toast} affiche des messages temporaires (succès, erreur).
\end{itemize}

\subsection{Gestion du thème sombre/clair}

Le thème est géré via la classe CSS \texttt{dark} sur l'élément \texttt{<html>} :

\begin{lstlisting}[language=Java, caption={Gestion du thème — extrait de app.js}]
isDark() {
    return document.documentElement
        .classList.contains('dark');
},

toggleTheme() {
    const html = document.documentElement;
    html.classList.toggle('dark');
    localStorage.setItem('theme',
        html.classList.contains('dark')
            ? 'dark' : 'light');
},
\end{lstlisting}

Le choix du thème est persisté dans \texttt{localStorage} et restauré au chargement de la page via un script inline dans \texttt{index.html}.

\subsection{Pages de l'application}

\begin{table}[H]
    \centering
    \begin{tabularx}{\textwidth}{l X}
        \toprule
        \textbf{Page} & \textbf{Description} \\
        \midrule
        Connexion & Formulaire pseudo/email + mot de passe, lien vers inscription \\
        Inscription & Formulaire pseudo + email + mot de passe, validation côté client \\
        Dashboard & Liste des capsules de l'utilisateur, badge d'invitations, bouton de création \\
        Détail capsule & Informations, liste des membres, souvenirs (si ouverte), formulaire d'invitation \\
        Créer capsule & Formulaire titre + description + date d'ouverture \\
        \bottomrule
    \end{tabularx}
    \caption{Pages de l'application SPA}
\end{table}


\section{Routage Apache}

Le fichier \texttt{.htaccess} assure le routage SPA : toutes les requêtes qui ne correspondent pas à un fichier ou dossier existant sont redirigées vers \texttt{index.html}.

\begin{lstlisting}[caption={Configuration .htaccess pour le routage SPA}]
RewriteEngine On
RewriteBase /testminiprjt/

# Ne pas rediriger les fichiers et dossiers existants
RewriteCond %{REQUEST_FILENAME} !-f
RewriteCond %{REQUEST_FILENAME} !-d

# Rediriger tout vers index.html
RewriteRule ^(?!api/)(.*)$ index.html [L,QSA]
\end{lstlisting}

\section{Design et interface utilisateur}

L'interface adopte un design \textbf{minimaliste} avec une palette de couleurs :
\begin{itemize}[noitemsep]
    \item \textbf{Blanc / Noir} : couleurs principales (fond et texte)
    \item \textbf{Vert nature} (\texttt{\#16A34A} — Tailwind \texttt{green-600}) : couleur d'accent pour les boutons, liens et éléments interactifs
    \item \textbf{Neutres} : échelle de gris \texttt{neutral} de Tailwind pour les bordures, textes secondaires et arrière-plans
\end{itemize}

Le \textbf{mode sombre} inverse les fonds (noir profond \texttt{neutral-950}) et adapte les contrastes grâce à la stratégie \texttt{darkMode: 'class'} de Tailwind CSS. L'utilisateur peut basculer entre les deux modes via un bouton soleil/lune, et son choix est persisté dans le navigateur.


% ══════════════════════════════════════════════
%  CHAPITRE 4 : CONCLUSION
% ══════════════════════════════════════════════
\chapter{Conclusion}

\section{Bilan}

Le projet \textbf{ChronoVault} a permis de concevoir et développer une application web complète répondant à la problématique initiale : offrir aux utilisateurs un outil collaboratif pour créer et partager des capsules temporelles numériques.

Les objectifs fixés ont été atteints :
\begin{itemize}[noitemsep]
    \item Un système d'authentification sécurisé avec hachage bcrypt et sessions PHP.
    \item Une gestion complète des capsules temporelles avec verrouillage par date.
    \item Un mécanisme d'invitation par pseudo permettant la collaboration.
    \item L'upload et l'affichage de souvenirs multimédia (texte + images).
    \item Une interface moderne, responsive et agréable avec support du mode sombre.
    \item Une architecture SPA fluide sans rechargement de page.
\end{itemize}

\section{Compétences acquises}

Ce projet nous a permis de renforcer nos compétences dans :
\begin{itemize}[noitemsep]
    \item La conception UML (cas d'utilisation, classes, séquence).
    \item Le développement d'API REST en PHP natif avec PDO.
    \item La construction d'une SPA en JavaScript vanilla.
    \item L'utilisation de Tailwind CSS pour un design rapide et cohérent.
    \item La gestion de la sécurité web (injections SQL, validation, hachage).
    \item Le travail collaboratif et la gestion de projet.
\end{itemize}

\section{Perspectives d'amélioration}

Plusieurs améliorations pourraient être envisagées pour enrichir l'application :
\begin{itemize}[noitemsep]
    \item \textbf{Notifications en temps réel} : utilisation de WebSockets pour notifier les invitations et l'ouverture de capsules.
    \item \textbf{Support vidéo} : permettre l'ajout de vidéos en tant que souvenirs.
    \item \textbf{Application mobile} : développer un client mobile natif ou PWA.
    \item \textbf{Chiffrement de bout en bout} : chiffrer les souvenirs pour garantir la confidentialité.
    \item \textbf{Système de commentaires} : permettre aux membres de réagir aux souvenirs après l'ouverture.
    \item \textbf{Emails automatiques} : notifier les utilisateurs par email lors d'une invitation ou de l'ouverture d'une capsule.
    \item \textbf{Migration vers un framework} : envisager Laravel (PHP) ou React/Vue (JS) pour une meilleure maintenabilité.
\end{itemize}

\vspace{1cm}
\noindent\rule{\textwidth}{0.4pt}
\vspace{0.3cm}

\begin{center}
    \textit{Projet réalisé par \textbf{ELHARCHY Mostafa} et \textbf{EL MESKI Karim}.}
\end{center}

\end{document}
